\begin{table}[H]
\centering
\small
\caption{Пример JSON-ответа модели после генерации комментария}
\label{tab:model-output-example}
\begin{tabular}{p{0.23\textwidth}p{0.67\textwidth}}
\toprule
Поле & Значение \\
\midrule
\texttt{t\_start} & \texttt{00:00:07.397} \\
\texttt{t\_end} & \texttt{00:00:07.871} \\
\texttt{commentary} & Хавертц делает пас низом после начального удара, но МакГрегор не может восстановить владение. \\
\bottomrule
\end{tabular}
\end{table}

В полном JSONL-файле этот ответ хранится вместе с технической информацией о стратегии группировки, номере payload, модели и статусе парсинга. Для сборки аудиодорожки используются прежде всего три поля: начало реплики, конец связанного эпизода и сам текст комментария.
